\section{پاسخ سوال ۴}
\begin{stepbox}
\field{افزایش کنتراست هیستوگرام}{
با توجه به شکل هیستوگرام، مقادیر پیکسل‌ها در بازه $r_{min} = 32$ تا $r_{max} = 160$ توزیع شده‌اند و قله (میانگین توزیع) دقیقاً در نقطه $r_{mid} = 96$ قرار دارد.

\textbf{الف) افزایش کنتراست به صورت خطی (Linear Contrast Stretching):}
در این روش، مقادیر پیکسل از بازه فعلی به کل بازه ممکن (۰ تا ۲۵۵) به صورت خطی نگاشت (Stretch) می‌شوند:
\[ s = \frac{r - r_{min}}{r_{max} - r_{min}} \times 255 = \frac{r - 32}{160 - 32} \times 255 = \frac{r - 32}{128} \times 255 \]
برای پیکسلی با مقدار ۹۶:
\[ s = \frac{96 - 32}{128} \times 255 = \frac{64}{128} \times 255 = 0.5 \times 255 = 127.5 \approx \textbf{128} \]

\textbf{ب) افزایش کنتراست به صورت Gamma Correction:}
فرمول کلی تصحیح گاما به شکل $s = 255 \times \left(\frac{r}{255}\right)^\gamma$ است.
برای بیشینه کردن کنتراست بصری (متعادل‌سازی هیستوگرام تاریک)، گاما به گونه‌ای انتخاب می‌شود که نقطه تمرکز هیستوگرام (یعنی ۹۶) به وسط بازه روشنایی (۱۲۸) منتقل شود. 
\[ 128 = 255 \times \left(\frac{96}{255}\right)^\gamma \implies \frac{128}{255} = \left(\frac{96}{255}\right)^\gamma \]
\[ \gamma = \frac{\ln(128/255)}{\ln(96/255)} \approx 0.705 \]
با این مقدار گاما، پیکسلی با مقدار اولیه ۹۶ به طور دقیق به مقدار \textbf{۱۲۸} نگاشت می‌شود تا کنتراست در میان‌تون‌ها بیشینه گردد.
}
\end{stepbox}
